% !TEX root = ../nc-cubic.tex
\subsection{一般の対角交換係数に対する有限位数制約}
\label{nds:sec:support}
\begin{proposition}\label{nds:prop:support}
$B=\C\langle z_1,\ldots,z_N\rangle/(z_i z_j-q_{ij} z_j z_i)$ で
$q_{ii}=1$, $q_{ji}=q_{ij}^{-1}\in\C^\times$ とし、
$f\ne0$ を中心三次式とする。$D^b(\operatorname{qgr} B/(f))$ が proper なら、
各 $i$ について $2e_i+e_{\phi(i)}$ が $f$ の PBW 指数に現れるように
$\phi(i)$ を選べる。従ってすべての $q_{ij}$ は有限位数である。
\end{proposition}
\begin{proof}
交換比を交代双指標 $b$ として $\mathbb Z^N$ 上に書く。
中心性と PBW 基底により $f$ の各指数は $b$ の radical にある。
それらの指数で割った群を $L$ とする。
$b$ は $L$ 上の双指標に降り、その交換子を持つ双線型乗数を選べる。
実際、有限生成可換群の巡回分解に関して $b$ の上三角成分を取ればよい。
この乗数による $L$ 次数付き twist により、可換多項式環は $B$ に、
同じ単項式支持を持つある可換三次 $F$ の商は $A=B/(f)$ に移る。
ここで通常の整数次数も保持する。$F$ の $L$ 次数は零である。

ある $i$ に必要な単項式が一つもないと仮定する。
$z_i$ が可逆な次数零 chart の可換モデルは
$D=\C[t_j:j\ne i]/(F(\ldots,1,\ldots))$, $t_j=x_j/x_i$ であり、
$L$ 次数は $e_j-e_i$ である。分母を消した式は非零で原点の極大
イデアルの二乗に入るので、原点は特異である。座標体 $k$ は
$\operatorname{Ext}_D^p(k,k)\ne0$ を任意に高い次数に持つ。
実際、原点の局所環での極小自由分解は正則性判定により無限に続く。
この Ext は局所化する前にも非零である。

非可換 chart $D'=(A[z_i^{-1}])_0$ は $D$ の同じ乗数による twist である。
斉次元の逆元をスカラーで補正して取れば Ore 局所化と twist は可換し、
両辺で通常次数零を取ればこの表示を得る。$k$ は原点の増大写像での加群である。
$L$ 次数付き有限自由加群による $k$ の解消を一段ずつ取れる
（可換側は noetherian であり核は斉次）。各項が有限階数なので、
$k$ に値を取る通常の Hom 複体は次数付き Hom 複体の全 $L$-shift にわたる直和である。
twist は次数付き加群の完全同値で自由加群の shift を自由加群の shift に移す。
任意の $k(\lambda)$ の twist は同じ次数に集中する同じ原点加群である。
従って、この直和のコホモロジーの次元は twist の前後で一致し、
$\operatorname{Ext}_{D'}^p(k,k)$ も任意に高い次数で非零となる。

これが射影圏でも同じ点の自己 Ext を与えることを確認する。
全次数付き加群の Ore 局所化は完全であり、torsion 加群を消すので
$\operatorname{QGr} A$ から局所化環の次数付き加群への完全関手となる。
その右随伴は制限の後に商関手を施す完全忠実充満関手である。
実際、局所化加群は torsion 加群からの Hom と一次 Ext が消える。
これは局所化の完全性と、制限が単射分解を保つことから従う。
局所化環の次数付き加群は、可逆な次数一の元 $z_i$ により
$D'$ 加群と同値である。原点加群を右随伴で戻した対象は
点加群 $A/(z_j:j\ne i)$ の像である。負次数を加えて $z_i$ を
可逆にすることで生じる余核は torsion だからである。
完全随伴により点の自己 Ext は上記の chart の自己 Ext と一致する。
局所 noetherian な $\operatorname{QGr} A$ では noetherian 対象間の
Yoneda Ext は $\operatorname{qgr} A$ でも同じである
（有限拡大列は端から順に noetherian 部分対象に縮められる）。
従って全 coherent 圏の proper 性に反する。

各行が単位ベクトル $e_{\phi(i)}$ の行列を $P_\phi$ とすれば、
$L$ は $2I+P_\phi$ の行で割った群 $L_\phi$ の商である。
$P_\phi$ の行和ノルムは一で、$2I+P_\phi$ は有理数上可逆である。
$L_\phi$ は有限であり $q_{ij}=b(e_i,e_j)$ は有限位数となる。
\end{proof}

\begin{corollary}\label{nds:cor:six}
$N=5$ では命題\ref{nds:prop:support}の仮定の下で $q_{ij}^6=1$ が全組で成立する。
\end{corollary}
\begin{proof}
$L_\phi$ の Smith 不変因子を $d_1\mid\cdots\mid d_5$ とすれば、
$\bigwedge^2 L_\phi$ の指数は $d_4$ である。
各 $\phi$ に対して $b$ はこの外積からの指標なので、
全交換係数の生成する群の指数は $d_4$ を割る。
$5^5$ 写像を変数置換で割って $2I+P_\phi$ の Smith 形を求めると
次表を得る。付属計算器
\texttt{verify\_diagonal\_quantum\_support.py} は軌道と全入力を構成する。
\[
\begin{array}{c|rrrrr}
d_4&1&2&3&4&6\\\hline
\text{軌道数}&5&14&16&1&11\\
\text{ラベル付き写像数}&504&1005&816&60&740
\end{array}
\]
$d_4=4$ の軌道は $0$ 始まりで $\phi=(0,0,0,1,2)$ である。
交換係数の指数が実際に四であるとする。
$W=L_\phi/\operatorname{rad}(b)$ において各座標の像を $w_j$ と書く。
$W$ の指数は四なので $3w_0=0$ から $w_0=0$、
$w_1=-2w_3$, $w_2=-2w_4$ となる。よって $w_3,w_4$ が生成し、
$b(w_3,w_4)$ は原始四乗根である。この二元との対合を取れば
$W=(\mathbb Z/4)^2$ の基底であることが従う。

有限群では中心商の全 coherent 圏は、可換三次 $Y=V(F)$ の
$W^*$-ねじれ同変 coherent 圏である。
実際、非退化双指標の有限ねじれ群環は中心がスカラーの半単純環で
一つの行列環となるので、定理\ref{nqt:thm:reduction}の射影生成族の
証明をそのまま適用できる。ここで $F$ は $W$ 次数零である。
proper 性から $Y$ が滑らかでなければならない。特異点があれば
その剰余体を全群へねじれ誘導し、忘却との随伴から高次自己 Ext を得るためである。
ところが $x_1,x_2$ の張る射影直線上では、どの三次単項式も重み零にならない。
この直線上の二次単項式の重みは $0$ または $w_1+w_2$ のみであり、
負の重みがそのどちらかとなる座標は $x_0$ だけである。
従って直線上で非零になり得る偏導関数は $\partial_0F$ のみである。
その二次式は複素射影直線上に零点を持つので特異点が生じ、矛盾する。
これで指数四を除き、表から結論する。
\end{proof}
